% SD4H Abstract — TARGET: 140-160 words
% Farquhar 5-sentence: (1) Achievement, (2) Why hard, (3) How, (4) Evidence, (5) Most remarkable
\begin{abstract}
Clinical predictors degrade across hospitals, but regulatory and hardware constraints often forbid changing the validated model.
We introduce \retr{}, a retrieval-guided input-space adapter that adapts the data, not the model, while the downstream predictor stays strictly frozen.
Most domain adaptation methods assume source labels are unavailable and the encoder is trainable. These assumptions fit many settings, but a complementary and under-studied regime arises when the target predictor's weights cannot be changed and source labels are available, a common situation in SaMD transfer, firmware-embedded models, and edge deployment. In this regime, alignment losses cannot propagate through a frozen encoder and classical DA is structurally blocked.
\retr{} encodes source windows, retrieves $k$ nearest target exemplars per timestep, and fuses them via cross-attention before decoding.
Across five ICU tasks (\eicu{}, \hirid{}~$\to$~\mimic{}) and five AdaTime benchmarks, a single method strictly outperforms every comparable baseline: all four frozen-backbone DA methods, an affine statistics baseline, end-to-end fine-tuning of the predictor, three end-to-end DA methods on the clinical classification tasks, and all eleven published end-to-end methods on AdaTime under that benchmark's own training protocol.
Adapted predictions exceed the target-native LSTM on two of five tasks and the source-native LSTM on four of five, with AKI reaching 91.1 \aucroc{} ($+16.1$ \aucpr{} over target-native).
\end{abstract}
